\scanpage{8}
\[
\begin{aligned}
&\deleted{\unsure{A|E|\operatorname{Ext}}}H^{r-i}(\widehat P,\Omega^{r}_{\deleted{X}P}(m))\\
&H^{\unsure{n|a}}(\widehat P,\Omega^{r}_{P}(m))\\
&\deleted{H^0(\widehat P\unsure{)|,)|,}\unsure{\mathord{-}|\mathord{=}}}\\
&\deleted{H^0(\widehat P,}
\end{aligned}
\]

\[
\resizebox{\linewidth}{!}{$\displaystyle\boxed{\begin{array}{cl}
\left.\left\{\begin{aligned}
&H^i(\widehat P,\Omega^r_{\widehat P}(m))=0 &&\text{si }0<i<n\\
&H^0(P,\Omega^r(m))\xrightarrow{\sim} H^0(\widehat P,\widehat\Omega^r(m))
\end{aligned}\right.\right|&
\Longleftrightarrow\ \begin{aligned}
&P-X\ \text{est de}\\[-2pt]
&\text{dim coh. }r-n\\[-2pt]
&[\text{et non plus}\\[-2pt]
&\quad\text{grand}]
\end{aligned}\\[3pt]
\begin{gathered}
\text{pour }m\text{ grand, ou tt }m\\
\text{(\unsure{au choix|en somme}) \unsure{on peut même}}\\[-2pt]
\text{remplacer }\Omega^r_P\text{ par }\underline O_P\ \unsure{\ldots|\text{ici}}
\end{gathered}&
\end{array}}$}
\]

Les schémas formels s'introduisent en Géométrie algébrique de la façon suivante :

1) Techniquement, pour énoncer et prouver de façon commode les deux th. fondamentaux des morphismes propres

% The source records the insertion at its destination. The box places its text
% below 3 as on the sheet; the arrow returns to the insertion point.
\begingroup
\newsavebox{\ScanEightItem}
\setbox\ScanEightItem=\vbox{\hsize=\linewidth\noindent
\tikz[remember picture,baseline=(sEightTwo.base)]{\node[draw,circle,inner sep=1pt] (sEightTwo) {2)};} \inserted{En théorie de Lefschetz - \unsure{Courant|Ce sont|Grauert}, \deleted{\unsure{pour|d'où}} \inserted{\unsure{le plus couramment des|le plus courant des}} \unsure{énoncés|énoncer} \deleted{les} \inserted{\unsure{donnant des|donnent des}} relations entre une variété (projective, ou locale) et son complété formel le long \unsure{``|d'}une section hyperplane''.}}

\noindent\tikz[remember picture,overlay]{\coordinate (sEightAt) at (.6,1.5ex);}3) En théorie formelle des Modules, pour \deleted{\unsure{pouvoir trouver|parfois trouver}} \inserted{\unsure{saisir}} une notion de familles modulaires\par
\smallskip\noindent\usebox{\ScanEightItem}\par
\noindent\tikz[remember picture,overlay]{\draw[->] (sEightTwo.west) .. controls +(-.6,.8) and +(-1.4,0) .. (sEightAt);}
\endgroup

4) En formalisme de dualité, les \unsure{duals|duels|énoncés} des groupes de cohomologie $H^p_Y(X,F)$ s'expriment par des $\operatorname{Ext}$ ou $H^{\unsure{*|q|p}}$ dans $\widehat X$ (complété formel le long de $Y$)

\margin{TSVP}
% archive sheet number: 7
